\section{Related Work}

Denoising diffusion probabilistic models (DDPMs) establish a likelihood-based generative
framework in which a neural network learns to iteratively reverse a fixed Gaussian noising
process applied to training data \cite{ho2020ddpm}. Since their introduction, diffusion models
have expanded considerably in both scale and architectural diversity, with transformer-based
backbones such as the Diffusion Transformer (DiT) demonstrating strong performance on
class-conditional image generation \cite{peebles2023dit}. The present work does not propose a
new diffusion architecture; rather, it examines practical fine-tuning behavior under fixed
training budgets, using established backbones as controlled experimental platforms.

Parameter-efficient fine-tuning (PEFT) methods reduce adaptation cost by confining gradient
updates to a small subset of model parameters. Early approaches introduced lightweight adapter
modules inserted between frozen transformer layers \cite{houlsby2019adapters}, while prompt-based
tuning methods condition frozen models through learned input representations
\cite{lester2021prompt}. LoRA offers a complementary formulation: rather than inserting new
modules, it augments frozen weight matrices with trainable low-rank factors that can be folded
back into the base weights at inference \cite{hu2022lora}. The rank parameter $r$ provides
direct, interpretable control over adaptation capacity, making LoRA a natural candidate for
systematic efficiency-quality trade-off analysis under fixed compute budgets.

A growing body of work has examined the limitations of static, uniform rank assignment in LoRA.
AdaLoRA \cite{zhang2023adalora} reparameterizes weight updates via singular value decomposition
and dynamically prunes less important singular values based on layer-wise importance scores,
allowing rank to adapt to the varying adaptation demands of different modules. Subsequent
methods have built on this direction: ElaLoRA \cite{chang2025elalora} enables both rank pruning
and expansion during training via gradient-derived importance estimates, while ARD-LoRA
\cite{ardlora2025} targets heterogeneous adaptation requirements across individual attention heads
rather than layers. Separately, rsLoRA \cite{kalajdzievski2023rslora} identifies a theoretical
limitation in the standard LoRA scaling factor, showing that the conventional $\alpha/r$
normalization causes gradient magnitude to collapse as rank increases, and proposes a corrected
$\alpha/\sqrt{r}$ scaling that stabilizes learning at higher ranks without additional inference
cost. More recently, timestep-dependent rank adaptation has been explored for diffusion model
personalization: T-LoRA \cite{soboleva2025tlora} demonstrates that higher diffusion timesteps
are disproportionately prone to overfitting, motivating rank schedules that vary with the
denoising trajectory.

While these adaptive and rank-aware methods represent important advances, they introduce
additional complexity---in the form of dynamic schedules, importance scoring, or modified
optimization regimes---that can obscure the baseline relationship between fixed rank and
adaptation quality. Our study deliberately occupies a complementary position: rather than
proposing a new rank-selection mechanism, we provide controlled empirical evidence of how
static rank choices interact with training budget under uniform optimization settings. This
evidence is directly relevant for practitioners who require reproducible baselines before
committing to more complex adaptive strategies.

Open-source frameworks have substantially lowered the barrier to reproducible diffusion
experimentation. Hugging Face Diffusers \cite{vonplaten2022diffusers} provides consistent
implementations of models and noise schedulers across a wide range of architectures, enabling
fair cross-study comparisons. Evaluation of generative quality commonly relies on the
Fr\'{e}chet Inception Distance (FID) \cite{heusel2017fid}, while CIFAR-10 remains a widely
used controlled benchmark that facilitates rapid and resource-efficient comparative studies
\cite{krizhevsky2009cifar}. The present study builds directly on these foundations, employing a
local \texttt{pytorch-fid} protocol with a fixed CIFAR-10 test reference to ensure evaluation
consistency across all experimental tracks.